\documentclass[11pt]{article}
\usepackage[margin=1in]{geometry}
\usepackage{booktabs}
\usepackage{longtable}
\usepackage{array}
\usepackage{hyperref}
\usepackage{xcolor}
\usepackage{titling}
\usepackage{parskip}
\setlength{\parindent}{0pt}

\hypersetup{
  colorlinks=true,
  linkcolor=blue!60!black,
  urlcolor=blue!60!black,
  citecolor=blue!60!black
}

\title{BVBRC-60 --- Independent Replication Report\\
\large \textit{Bacillus megaterium} NCT-2 Complete Genome (Wang et al., 2020)}
\author{Replication set: BVBRC-100 (rank 48, BVBRC\_TOPUP85\_2026-06-26)}
\date{}

\begin{document}
\maketitle

\section*{Paper}
Wang B, Zhang D, Chu S, Zhi Y, Liu X, Zhou P.
\emph{Genomic Analysis of Bacillus megaterium NCT-2 Reveals Its Genetic Basis for the Bioremediation of Secondary Salinization Soil.}
\emph{International Journal of Genomics}, 2020, Article 4109186.
\href{https://doi.org/10.1155/2020/4109186}{doi:10.1155/2020/4109186} $\cdot$
PMID 32190639 $\cdot$ PMCID PMC7066406 (open access, CC-BY).

\textbf{Target organism:} \emph{Bacillus megaterium} NCT-2 (reclassified \emph{Priestia megaterium} NCT-2), a soil-bioremediation / plant-growth-promoting isolate; CGMCC No.~4698.

\textbf{BV-BRC-mappable workflow:} Comprehensive Genome Analysis (assembly + RASTtk-style annotation) + Similar Genome Finder / PlasmidFinder + Phylogenetic Tree.

\section{Paper Summary}

NCT-2 is a nitrate-uptaking, salinity-adapted, phosphate-solubilizing \emph{B.\ megaterium} isolated from secondary-salinized greenhouse soil. The authors sequenced its complete genome with a hybrid Illumina HiSeq~4000 + PacBio RSII strategy, resolving a $\sim$5.19~Mb circular chromosome plus 10 indigenous plasmids (total 5.88~Mb, $\sim$37.87\% GC). They annotate the genome (GO/KEGG/COG/RAST), place it phylogenetically (CVTree + 16S NJ + MAUVE) closest to \emph{B.~megaterium} DSM~319 then QM~B1551, and mine gene inventories that explain the strain's bioremediation/PGPR phenotype: nitrogen metabolism (nitrate\,$\rightarrow$\,ammonia\,$\rightarrow$\,glutamate assimilation), phosphate solubilization/uptake, IAA (Trp-dependent, incomplete) biosynthesis, and osmotic/oxidative stress tolerance.

\section{Claims Table}

\begin{longtable}{@{}p{0.6cm}p{9cm}p{2.2cm}p{1.4cm}p{1.4cm}@{}}
\toprule
ID & Claim & Type & Testable? & Tested? \\
\midrule
\endhead
C1 & 1 circular chromosome ($\sim$5.19~Mb) + 10 plasmids; total 5.88~Mb & architecture & Yes & Yes \\
C2 & Whole GC 37.87\%; chromosome 38.2\%; plasmid GC 33.7--37.0\% & sequence stat & Yes & Yes \\
C3 & 6{,}039 genes; 5{,}606 CDS; 203 RNA; 230 pseudogenes; 142 tRNA; 53 rRNA (19$\times$5S, 17$\times$16S, 17$\times$23S) & annotation & Yes & Yes \\
C4 & Comparative Table~1: NCT-2 largest genome/most CDS among 6 \emph{Bacillus}; comparator sizes/GC & comparative & Yes & Yes \\
C5 & Phylogeny: NCT-2 most homologous to DSM~319, then QM~B1551 & phylogeny & Yes & Yes \\
C6 & Gene inventories: N-metabolism, phosphate (pstSCAB, alk.~phosphatase, glucose 1-DH), IAA (aldehyde DH + amidase), stress (glycine betaine ABC, betaine-aldehyde DH, SOD, catalase) & functional & Yes & Yes \\
C7 & Wet-lab isolation/provenance + HiSeq+PacBio hybrid workflow; CGMCC 4698 & provenance & No & No \\
\bottomrule
\end{longtable}

\section{Independent Method}

All data pulled from public NCBI; all reproduced numbers computed from downloaded files (Python~3 stdlib for FASTA/GFF parsing; fastANI for ANI). No paper numbers were copied into the computation.

\begin{enumerate}
\item \textbf{Paper}~---~Europe PMC full-text XML for PMC7066406, extracted deposited accessions (chromosome CP032527.2 + plasmids CP032528--CP032537).
\item \textbf{Study genome}~---~NCBI Datasets v2 REST: \texttt{GCA\_000334875.3} (ASM33487v3, Complete Genome) with \texttt{GENOME\_FASTA,GENOME\_GFF,PROT\_FASTA}. The 11 replicon accessions match the paper's Data Availability exactly. (RefSeq GCF\_000334875.3 v.3 replaced the earlier 204-contig draft v.1.)
\item \textbf{Genome stats}~---~per-replicon length and GC computed directly from FASTA (\path{report/evidence/genome_stats.json}).
\item \textbf{Annotation counts}~---~feature-type tally from \texttt{genomic.gff} (\path{report/evidence/annotation_counts.json}); protein count from \texttt{protein.faa}.
\item \textbf{Comparators}~---~NCBI Datasets REST for the five Table-1 reference genomes; sizes/GC computed identically (\path{report/evidence/comparative_genome_table.tsv}).
\item \textbf{Phylogeny}~---~\texttt{fastANI -q NCT-2 --rl <5 comparators>} (\path{report/evidence/ani_nct2_vs_comparators.tsv}).
\item \textbf{Functional genes}~---~grep of deposited protein-product annotations for each claimed inventory (\path{report/evidence/functional_genes_found.txt}).
\item \textbf{Verdict}~---~LLM judge (Argo \texttt{gpt-5.2}, free proxy) over the claim/result JSON (\path{report/evidence/llm_judge_verdict.txt}).
\end{enumerate}

Tool versions/sources: NCBI Datasets v2 REST (free, no auth); Europe PMC REST; fastANI (\path{/usr/local/bin/fastANI}); Python~3.

\section{Results vs Paper}

\subsection*{C1 --- Genome architecture (AGREE)}
\begin{tabular}{@{}lll@{}}
\toprule
 & Paper & Reproduced (GCA\_000334875.3) \\
\midrule
Chromosome & 5.19~Mb, circular & CP032527.2 = 5{,}193{,}616~bp \\
Plasmids & 10 (9{,}625~bp $\rightarrow$ >132~kb) & 10 (CP032528--CP032537; 9{,}625 $\rightarrow$ 132{,}087~bp) \\
Total & 5.88~Mb & 5{,}883{,}957~bp = 5.88~Mb \\
\bottomrule
\end{tabular}

\subsection*{C2 --- GC content (AGREE / MINOR-DIFF)}
\begin{tabular}{@{}lll@{}}
\toprule
 & Paper & Reproduced \\
\midrule
Whole genome & 37.87\% & 37.78\% \\
Chromosome & 38.2\% & 38.18\% \\
Plasmid range & 33.7--37.0\% & 33.65--37.02\% \\
\bottomrule
\end{tabular}

\subsection*{C3 --- Annotation totals (AGREE / MINOR-DIFF)}
\begin{tabular}{@{}lll@{}}
\toprule
Metric & Paper & Reproduced \\
\midrule
Genes (incl.~pseudo) & 6{,}039 & 6{,}038 \\
CDS / proteins & 5{,}606 & 5{,}605 proteins (5{,}845 CDS features) \\
RNA genes & 203 & 203 \\
Pseudogenes & 230 & 230 \\
tRNA & 142 & 142 \\
rRNA (5S,16S,23S) & 53 (19,17,17) & 53 (19,17,17) \\
\bottomrule
\end{tabular}

\subsection*{C4 --- Comparative table (AGREE)}
\begin{tabular}{@{}lll@{}}
\toprule
Strain & Paper size / GC & Reproduced size / GC \\
\midrule
NCT-2 & 5.88 / 37.8 & 5.88 / 37.78 \\
QM~B1551 & 5.52 / 37.97 & 5.52 / 37.93 \\
DSM~319 & 5.10 / 38.1 & 5.10 / 38.13 \\
\emph{B.~subtilis} 168 & 4.22 / 43.5 & 4.22 / 43.51 \\
\emph{B.~cereus} Q1 & 5.51 / 35.5 & 5.51 / 35.47 \\
\emph{B.~licheniformis} DSM~13 & 4.22 / 46.2 & 4.22 / 46.19 \\
\bottomrule
\end{tabular}

NCT-2 has the largest genome, as claimed.

\subsection*{C5 --- Phylogeny (AGREE)}
fastANI, NCT-2 vs comparators:
\begin{itemize}
\item \textbf{DSM~319 $\rightarrow$ 98.2\%} (closest) \checkmark
\item \textbf{QM~B1551 $\rightarrow$ 96.5\%} (second) \checkmark
\item \emph{B.~subtilis} 168 / \emph{B.~cereus} Q1 / \emph{B.~licheniformis} DSM~13 $\rightarrow$ below fastANI's $\sim$80\% reporting cutoff (distant, different species) \checkmark
\end{itemize}
Ordering exactly matches the paper ("most homologous to DSM~319 and then QM~B1551").

\subsection*{C6 --- Functional gene inventories (AGREE)}
All claimed inventories present in the deposited annotation (see \path{functional_genes_found.txt}):
\begin{itemize}
\item \textbf{Nitrogen:} NarK/NasA nitrate transporter, nitrite reductase NirD, nitroreductases, NifU, P-II nitrogen regulator, ammonium transporter, glutamate synthase large/small (GOGAT), type-I glutamate--ammonia ligase (GS), formate/nitrite transporter.
\item \textbf{Phosphate:} alkaline phosphatase, glucose 1-dehydrogenase, PstS/PstA/PstC + phosphate ABC ATP-binding (pstSCAB).
\item \textbf{IAA:} aldehyde dehydrogenase + amidase (consistent with the paper's "incomplete Trp-dependent pathway").
\item \textbf{Stress/osmoadaptation:} glycine betaine ABC transporter (opu-type), betaine-aldehyde dehydrogenase (gbsA), superoxide dismutase, catalase.
\end{itemize}

\subsection*{C7 --- Provenance/method (NOT-TESTABLE)}
Isolation, CGMCC deposit, and HiSeq+PacBio hybrid workflow cannot be re-verified from deposition alone; "complete genome" status is consistent with the NCBI record.

\section{Verdict}

\textbf{Coverage:} 6/6 testable claims tested (1.00). \quad
\textbf{Agreement:} 6/6 AGREE or MINOR-DIFF (1.00). \quad
\textbf{LLM judge} (Argo gpt-5.2): \textcolor{green!50!black}{\textbf{REPLICATED}}.

Every genome-derived claim reproduces on the independently downloaded deposited genome with exact or sub-0.3\% rounding-level differences. The only non-matching element is the inherently non-testable wet-lab provenance.

\begin{center}
\fbox{\textbf{Verdict: REPLICATED}}
\end{center}

\section{Genuine Critique}

This section states what the replication \emph{did} and did not actually settle, and where a reader should keep their prior loose.

\subsection*{What was genuinely tested}
The replication re-computed, from the deposited assembly and its GFF/proteome, every quantitative genome-scale claim in the paper: replicon count and sizes, per-replicon and whole-genome GC, gene / CDS / RNA / pseudogene / tRNA / rRNA counts, the 6-strain comparative size/GC table, and phylogenetic rank-ordering via ANI. All six testable claims agreed within a few 0.01\%--0.3\% rounding-level offsets attributable to differences between the authors' original count and the current NCBI v.3 assembly annotation (one gene, one CDS, 0.09\% whole-genome GC). Functional inventories were confirmed by keyword grep of protein products against the deposited annotation. This is a substantive and near-perfect reproduction of the \emph{informatic} content of the paper.

\subsection*{What the replication does \textbf{not} show}
\begin{itemize}
\item \textbf{Wet-lab provenance is unverified (C7).} The claim that this sequence corresponds to a real isolate CGMCC~4698 obtained from a specific greenhouse soil, and that it was produced by the stated HiSeq~4000 + PacBio~RSII hybrid strategy, cannot be checked from the deposition. A misassignment at deposit time would not be caught by any test performed here.
\item \textbf{Phenotype $\rightarrow$ genotype attribution is not tested.} The paper's central narrative is that the observed genes \emph{explain} NCT-2's salinization/PGPR phenotype. The replication only confirms that the annotated genes are present; it does not test that they are expressed, that they function as annotated in this strain, that their activity accounts for the measured phenotype, or that removing them abolishes it. Presence of \emph{opu}, \emph{gbsA}, SOD, catalase in a \emph{Bacillus} genome is unsurprising and shared with many non-remediation strains.
\item \textbf{Homology by name, not by structure.} The functional gene evidence is grep-over-product-strings, not HMM or orthology assignment. A miscurated product line would score identically to a genuine hit. In particular the "incomplete Trp-dependent IAA pathway" claim is supported only by observing an aldehyde dehydrogenase and an amidase, both of which are extremely broad-specificity enzyme families.
\item \textbf{Phylogenetic method substitution.} The paper used CVTree + 16S NJ + MAUVE; the replication used fastANI. The \emph{rank ordering} of the two \emph{B.~megaterium} comparators matches (DSM~319 closest, then QM~B1551), which is the substantive claim, but the two methodologies measure related-but-different quantities. A more faithful replication would rerun CVTree; that is a scope choice, not a defect, but it should be named.
\item \textbf{Comparator set is inherited, not chosen.} Table~1 uses the same six strains as the paper. The replication does not test whether NCT-2's "largest genome" claim survives against a broader \emph{Priestia}/\emph{Bacillus megaterium} pangenome, which now contains many complete genomes larger than QM~B1551.
\item \textbf{Assembly-version drift.} The reproduced numbers come from GCA\_000334875.3 (v.3); the paper used the original v.1 deposit. The 1-gene / 1-CDS delta almost certainly reflects annotation-pipeline evolution rather than sequence change, but this is a plausible interpretation, not a proven one.
\item \textbf{LLM judge is corroborative, not adjudicative.} The final REPLICATED label is emitted by a language-model judge over the numeric agreement table; the judge does not add evidence, only summarizes.
\end{itemize}

\subsection*{Bottom line of the critique}
The paper's \emph{informatic} description of the NCT-2 genome is real and reproducible; a reader can trust the size, structure, annotation counts, comparative table, and phylogenetic ordering. The paper's \emph{functional / adaptive} story --- that this specific gene set explains this specific bioremediation phenotype --- is \textbf{consistent with but not established by} the deposited data, and remains an open experimental question. The REPLICATED verdict applies to the former, not the latter.

\end{document}
